\documentclass[11pt]{article}
\usepackage[utf8]{inputenc}
\usepackage[T1]{fontenc}
\usepackage{lmodern}
\usepackage[margin=1in]{geometry}
\usepackage{csquotes}
\usepackage[style=authoryear,backend=biber,sorting=nyt,maxbibnames=99,giveninits=true,uniquename=init]{biblatex}
\usepackage[hidelinks]{hyperref}
\addbibresource{references.bib}
\begin{document}
\section*{Tile Drainage Influence on Streamflow Dynamics}

\subsection{Overview}

Agricultural subsurface (tile) drainage exerts a pervasive and quantifiable influence on streamflow, reshaping how water is partitioned among pathways and how the resulting hydrograph behaves across both event and non-event periods. The evidence base assembled here spans edge-of-field and plot monitoring, isotopic and geochemical tracer studies, hydrograph-separation and recession analyses, physically based and analytical modeling, regional data synthesis, and a global meta-analysis, and it converges on several well-supported conclusions. First, tile drainage constitutes a major component of watershed water yield, frequently supplying on the order of 40--65\% of streamflow during much of the year, with the largest contributions during spring and other wet, low-intensity periods \parencite{King2014,Thomas2016,Schilling2019}. Second, drainage systematically raises low flows and the baseflow fraction, dampens intermediate flows, and has comparatively little effect on the very largest floods \parencite{Sloan2017,Schilling2008,Rathore2025}. Third, artificial drainage \emph{homogenizes} landscape hydrologic response, shortening and stabilizing travel times and rendering hydrologic signatures more scale-invariant, with the distinction between tiled and non-tiled behavior fading above characteristic catchment sizes of roughly 2000~km$^2$ \parencite{BolandBrien2014,DaneshYazdi2016,FoufoulaGeorgiou2015}. Fourth, runoff generation in drained fields is governed by pathway-activation thresholds tied to antecedent moisture and rainfall, with tile flow typically the dominant and most persistent export pathway \parencite{Williams2019,Williams2018,Macrae2019}. Finally, effects on flashiness, peak flows, and streamflow drought are conditional on slope, soil, storm size, and drainage intensity, so drainage can either enhance or attenuate event response depending on setting \parencite{Zhou2025,Rathore2025,Adelsperger2024}. The subsections below synthesize these lines of evidence in turn, together with the mechanistic, tracer, and modeling work that underpins them.

\subsection{Tile drainage contributions to streamflow volume}

\subsubsection{Annual and event-scale partitioning}

A consistent body of watershed-scale work establishes that tile drains route a large and measurable share of precipitation to streams. In an Ohio watershed, 21\% of precipitation (206~mm) was recovered annually through tile drainage, which accounted for 41\% of the annual and 47\% of the monthly total watershed discharge; annual discharge averaged 498~mm, or 49\% of precipitation, split into 41\% baseflow and 59\% storm-event discharge \parencite{King2014}. Equivalent-medium modeling of an Iowa watershed in 2014 attributed 97~mm of tile drainage at the outlet---13\% of precipitation and 39\% of streamflow---within an annual water balance of $P=770$~mm, $ET=522$~mm, and $Q=248$~mm ($Q/P=0.32$) \parencite{Thomas2016}. In the Boone River (Iowa), end-member mixing analysis (EMMA) and SWAT bracketed the tile contribution at roughly 46\% to 54\% of discharge, with SWAT estimating about 53.7\% of Boone River and 65.8\% of Lyons Creek annual water yield over a 30-year simulation, and EMMA yielding 31--53\% (maximum) or 19--37\% (minimum) of streamflow over six years \parencite{Schilling2019}. Snowmelt-dominated SWAT modeling in a US--Canada watershed found tile flow equal to about 16\% of annual precipitation, or roughly 37\% of water yield, while tiling reduced surface runoff and increased subsurface flow, water yield, and drainage \parencite{Rahman2014}. Field-scale export can be even more tile-dominated: in Ohio, tile discharge accounted for $83\pm20$\% of water export where both surface runoff and tile flow were measured \parencite{Williams2018}, and most discharge from an Indiana closed depression was delivered by subsurface tile drains rather than surface flow \parencite{Williams2019}.

The partitioning between tile and total flow varies systematically with hydrologic condition and drainage area. In the Iowa equivalent-medium study, tile drainage represented on average 30\% of streamflow during precipitation events but 61\% during non-event periods; event-period $Q_T/Q$ increased with drainage area (mean slope $0.10$, mean $R^2=0.83$) whereas non-event $Q_T/Q$ decreased with area (mean slope $-0.17$, mean $R^2=0.82$), and the annual $Q_T/Q$ rose log-linearly from 0.29 at 6~km$^2$ to 0.39 at 45~km$^2$ ($R^2=0.98$) \parencite{Thomas2016}. Tile-flow magnitude also increased with contributing area, with full-pipe flow more frequent in smaller edge-of-field tiles than in larger mains \parencite{King2014}. The equivalent saturated hydraulic conductivity representing this behavior scaled with drainage density, from 164~m~d$^{-1}$ for collector drainage to 530~m~d$^{-1}$ at 10~m spacing ($R^2=0.99$) \parencite{Thomas2016}.

\subsubsection{Seasonal dynamics}

Tile contributions are strongly seasonal, peaking when soils are wet and evapotranspiration is low. Tile discharge in Ohio was seasonally variable, with median monthly summer discharge (0.9~mm) far below winter (23.4~mm), spring (10.2~mm), and fall (15.6~mm) values \parencite{King2014}. In the Boone River, tile drainage dominated flow from March through June, and October--November yields were largely groundwater plus tile discharge \parencite{Schilling2019}. DRAINMOD simulations in Iowa indicated the largest drainage effect on baseflow from March to July---more than 40\% more baseflow in drained watersheds and about 60\% more in April, May, and June \parencite{Schilling2008}. Expanded subsurface drainage was projected to decrease winter monthly flows, increase late-summer and fall flows, and have only a negligible effect on the following spring's flow volumes \parencite{Rahman2014}. In cold regions, tile flow generally occurred during snowmelt, tracking the synchrony between snow-water-equivalent depletion and drainage, with a seasonal absence of tile flow from May to September \parencite{Kompanizare2024}. Tile outflow can be a small fraction of rainfall during growing seasons except in continually wet years, reflecting how water-table position relative to tile depth gates drainage: in Iowa, mean water tables in a well-drained block lay below tile depth 72\% of the time and gradients steepened during wet periods, promoting lateral flow \parencite{Logsdon2020}.

\subsection{Modification of the flow regime}

\subsubsection{Low flows and baseflow}

A robust and repeated finding is that tile drainage elevates low flows and the baseflow fraction. Modeling and monitoring in Iowa showed drained landscapes carried far higher baseflow than undrained: average annual streamflow of 208~mm with 150~mm baseflow (72\%) at 24~m tile spacing versus 201~mm streamflow with 52~mm baseflow (25\%) with no tiles; tile drainage was concluded to primarily affect the baseflow portion of the hydrograph, though a detectable stormflow signal also appeared in tile hydrographs \parencite{Schilling2008}. Ohio simulations similarly found that tiles greatly increased low flow while decreasing higher flows to a lesser extent, raised annual average daily flow by about 5\% (via reduced evapotranspiration in the drier tiled soils), and lowered the probability of the water table entering the root zone ($<0.5$~m) from about 50\% without tiles to about 5\% with tiles \parencite{Rathore2025}. Routing analysis reinforced the general pattern that tile drains increase low flows, decrease intermediate or medium flows, and have minimal effect on the largest flows \parencite{Sloan2017}. In Latvia, drainage installation raised monthly mean and minimum river discharge, with significant increases in minimum monthly discharge in April, June, July, and August; mean runoff increased by 132--186\% during summer--autumn low-water periods (June--October) and by 53--286\% in the winter low-water period, while annual mean, minimum, and maximum discharge rose by 46\%, 158\%, and 26\%, respectively \parencite{Apsite2017}. Channelized (drained) headwater streams in Ohio likewise showed greater minimum and 10th-percentile discharge and far fewer zero-discharge days (8~d~yr$^{-1}$) than natural streams (87.5~d~yr$^{-1}$), inverting part of the accepted headwater paradigm \parencite{King2009}.

\subsubsection{Peak flows and flood response}

Effects on peak flows are consistently modest for the largest events but scale- and condition-dependent for smaller ones. Routing modeling found peakflow reductions concentrated at intermediate scales (100--1{,}000~km$^2$) and smaller at the largest scales ($>10{,}000$~km$^2$); for one event, reduction fell from 26\% at Wolf Creek watershed to 11\% at the Cedar River Basin outlet, and for another from 18\% at field scale to 10\% at the basin outlet, with negligible field-scale effect on the largest floods \parencite{Sloan2017}. The spatial placement of tiled hillslopes mattered: two watersheds with identical drained fractions could respond differently, and draining headwaters produced the greatest peak reduction because the width function clustered hillslopes there \parencite{Sloan2017}. Field-scale drained and undrained scenarios in Iowa showed no statistically significant difference in mean annual peak discharge for any of six soil types, indicating that drainage chiefly alters the flow pathway rather than the largest peaks \parencite{Sloan2016}. In Ohio ATS simulations, tiles could either enhance or attenuate event peaks depending on event size and antecedent conditions---enhancement for sub-threshold or dry-condition events, reduction for larger events---and streamflow was reduced by drainage in 23 of the 25 largest events, with a decision boundary separating positive from negative peak reduction 75\% of the time \parencite{Rathore2025}. Under expanded tiling, smaller peaks became more frequent while larger peaks declined (turning point at a 3.44-yr return period, not statistically significant), and steeper rising limbs with slightly shorter time-to-peak implied that reduced large-event peaks arose from reduced flow volumes rather than altered timing \parencite{Rahman2014}. Distributed hillslope modeling similarly found that drainage tends to decrease peak flow while increasing recession volume, reproducing observed behavior even with only 21\% of the area tiled \parencite{Fonley2021}.

\subsubsection{Flashiness}

Flashiness responds to drainage in a soil- and slope-dependent manner rather than monotonically. A global meta-analysis of 464 watersheds in 13 countries found that stream flashiness (Richards--Baker Index, RBI) was similar overall regardless of tile-drainage extent, but diverged by setting: in flat ($<2$\% slope) watersheds, medium and high drainage extents had higher RBI than low extent, whereas in sloping watersheds low extent had the highest RBI; in low-infiltration soils RBI was higher in lightly drained watersheds, while in high-infiltration soils RBI was similar across classes but the baseflow index (BFI) was higher where drainage was low \parencite{Zhou2025}. Drainage extent itself was positively associated with agricultural land use and clay content and negatively with slope and depth to water table \parencite{Zhou2025}. Field-scale Iowa simulations showed drainage decreased the flashiness index for low-conductivity soils and increased it for high-conductivity soils, thereby homogenizing response across soil types, with a nonlinear, U-shaped dependence on drain spacing and an optimal spacing near 15~m for the Webster soil \parencite{Sloan2016}. Ohio modeling found overall flashiness decreased with tiles, with a nonmonotonic minimum near 50~m spacing \parencite{Rathore2025}. Across 30 Nordic and Baltic catchments, the flashiness index was negatively related to BFI and positively related to the share of agricultural area under artificial drainage and to drainage intensity; computing flashiness from hourly rather than daily discharge increased it by nearly a factor of three on average (range 1.2--4.8) \parencite{Deelstra2014,Deelstra2014b}.

\subsubsection{Streamflow drought and long-term trends}

At the regime scale, tile drainage reshapes streamflow-drought characteristics. Across US watersheds, drainage significantly increased streamflow drought duration and intensity while reducing drought frequency, with effects comparable in magnitude to water-table depth and precipitation seasonality; each 10\% increase in tile-drained area lengthened drought duration by 0.03~d, raised intensity by 12\%, and lowered frequency by 0.10 events~yr$^{-1}$, and the result was robust to $\pm25$\% variation in the drained-area estimate \parencite{Adelsperger2024}. Regional trend analysis in Indiana found increasing low- and medium-flow trends, whereas high-flow trends were less linked to precipitation in heavily drained northern Indiana; accounting for autocorrelation and long-term persistence sharply reduced the count of significant low- and medium-flow trends, and the weakened precipitation linkage in drained areas suggested a drainage influence on streamflow trends warranting further study \parencite{Kumar2009}. A synthesis of seven Midwestern rivers concluded that inter-annual variability in streamflow, baseflow, and nitrogen load was controlled primarily by current- and previous-year precipitation rather than land cover or fertilizer, with adjusted $R^2$ up to 0.77 for baseflow in the Minnesota River, implying that recent river N-load increases are largely climate-driven and that tile-drained water will require remediation to reduce loads \parencite{Baeumler2020}.

\subsection{Scale dependence and hydrologic homogenization}

A central, well-supported theme is that artificial drainage homogenizes catchment response and reduces the scale dependence of hydrologic signatures. Data synthesis across Iowa showed that non-tiled watersheds exhibit decreasing area-normalized peak discharge and increasing baseflow ratio with scale, whereas tiled basins behave scale-invariantly; within-basin variability in recession coefficients declines as drained fraction rises, flow-duration curves show greater variability in non-tiled than tiled regions, and the tiled/non-tiled distinction disappears above roughly 2200~km$^2$, even though mean annual runoff ratio does not vary significantly with tiling \parencite{BolandBrien2014}. Travel-time analysis in Minnesota's Redwood basin found time-averaged mean travel time fell 41\%---from 77~d before land-use change to 45~d after, with the largest reduction in the growing season---while year-to-year variability was substantially reduced and travel time became independent of catchment size above about 200~km$^2$, indicating a more homogenized response under extensive tiling \parencite{DaneshYazdi2016}. Companion analysis showed streamflow dynamics became more linear after 1976 even though precipitation dynamics changed little, with strengthened dependence of streamflow on previous-day rain and more punctuated falling limbs \parencite{FoufoulaGeorgiou2015}. Isotope and transit-time work in a managed UK catchment found short mean transit times ($<1$~yr) with even faster artificial drains, and concluded that management---particularly artificial drainage---homogenized distributed catchment responses despite variation in subcatchment characteristics \parencite{DimitrovaPetrova2020}. In an Indiana tile-drained watershed, subwatersheds shared very similar spatial distributions of available water storage and hydrograph shape, so the basin behaved as an effectively homogeneous unit describable by a single exponential travel-time distribution with a 2--3~d mean \parencite{Basu2010}. Even so, tile contributions to annual discharge remained detectable as basins grew---beyond 16{,}000~km$^2$ in the Boone River---while diminishing at the largest scales \parencite{Schilling2019}.

\subsection{Runoff generation, flow pathways, and thresholds}

\subsubsection{Pathway activation and thresholds}

Runoff generation in drained fields is governed by thresholds linking antecedent moisture and rainfall to pathway activation. In an Indiana closed depression monitored over 289 rainfall events, rainfall amount alone related positively and linearly to cumulative flow with no visible threshold, but a clear threshold emerged when antecedent soil-moisture deficit (SMD) and rainfall were summed; thresholds varied by pathway and landscape position and were nearer zero at shoulder-slope than footslope sensors \parencite{Williams2019}. Subsurface tile flow generally began earlier, peaked earlier, and lasted longer than surface flow, and connected localized zones of saturation to enable rapid connectivity \parencite{Williams2019}. Isotopic work at the same site showed tile-drain celerity ($Q/P$) was largely controlled by antecedent wetness, exhibiting threshold behavior in SMD-plus-rainfall and in maximum groundwater depth relative to tile depth, with preferential flow more important under dry conditions and stored pre-event water dominating when wetter \parencite{Williams2021}. In near-level vertisolic clay in Canada, overland flow generally activated before groundwater and, in the tiled field, before tile flow in 14 of 18 events; tile-flow activation time related significantly to maximum rainfall intensity ($r^2=0.68$) and antecedent precipitation ($r^2=0.57$), improving to $r^2=0.78$ combined, and direct preferential flow into tiles was uncommon, implying a largely top-down sequence and tiles somewhat decoupled from the surface \parencite{Kokulan2019}.

Across US and Canadian sites, surface runoff was less frequent than tile flow: of 795 events, 43\% produced both surface and tile flow and 57\% produced tile flow only, and among dual-pathway events surface flow contributed a median of 12\% of total flow (range $<$1--97\%) \parencite{Macrae2019}. Surface-runoff frequency varied by texture and region (34--54\% of events on clay; 24--32\% at Ontario silt-loam sites; 48--59\% at Ohio silt-loam sites; 29\% at the clay-loam site), and event-response types (Tile, Surface, Simultaneous, Saturation) partitioned systematically by soil texture and season \parencite{Macrae2019}. In Kansas tile-outlet-terrace systems, baseflow contributed 38\% of discharge at one site, 51\% at another, and 46\% at a third, with short events dominated by quickflow and long events carrying larger baseflow components; antecedent conditions and storm size were stronger drivers of hydrologic and biogeochemical response than tile-inlet density alone \parencite{Stops2022}.

\subsubsection{Preferential versus matrix flow}

Tracer and geochemical studies partition tile flow into fast preferential and slower matrix components whose balance shifts with season and wetness. In a French field, seasonal runoff generation differed: under dry conditions rainwater transferred directly to tile drains through macropores, whereas under saturated conditions pre-event soil water dominated and rainfall contributed up to 25\% of tile-drain water; surface and tile drainage transferred 44\% and 20\% of annual effective rainfall, respectively, and 65\% of infiltrating water was not captured by tiles \parencite{Gaillot2023}. Catchment-scale monitoring there identified two soil-column transfer pathways---slow drainage of saturated soils and fast preferential flow through macropores/cracks---with the fast peak contributing about 15\% of water and sediment fluxes on average \parencite{Grangeon2021}. In Ohio, master-recession analysis distinguished two reservoirs (recession coefficients 0.9 and 0.25~d$^{-1}$; Reservoir~1 accounting for 54\% of subsurface flow), and a coupled pathway-connectivity separation apportioned annual tile discharge into 9\% quick-old, 39\% quick-new, 42\% slow-old, and 10\% slow-new, showing that new water routes through both quick- and slow-flow reservoirs and that groundwater recharge contributes importantly to tile fluxes \parencite{Nazari2021}. In Canada, preferential flow estimated from temperature and electrical conductivity agreed strongly ($r=0.83$), accounted for a higher fraction of tile flow in clay than loam soil, and correlated with phosphorus loads, indicating a disproportionate role in P export \parencite{Pluer2020}. Swedish drain-outflow modeling attributed autumn flow events to preferential flow and spring events to matrix flow under saturated conditions, with DRAINMOD-predicted cumulative infiltration explaining most variation in phosphorus concentrations \parencite{Salazar2011}. Isotopic hydrograph separation in Indiana showed groundwater was the primary source to the tile drain, followed by soil water, with minimal direct precipitation contribution and mean transit time increasing with depth (12--20~d at 10~cm to 245~d for tile drainage) \parencite{Williams2021}.

\subsubsection{Temporal concentration of flow and export}

Tile and surface flow are highly intermittent and temporally concentrated. In Ohio fields, surface runoff occurred on average $41\pm42$~d~yr$^{-1}$ and tile drainage $240\pm69$~d~yr$^{-1}$; 80\% of annual surface runoff occurred in 2.5\% of the year (9~d) and 80\% of tile discharge in 21\% (77~d), with Gini coefficients of 0.97--0.98 (surface) and 0.79--0.88 (tile) including zero-flow days \parencite{Williams2018}. Nutrient loads were similarly concentrated, and flow-duration/Lorenz analysis showed a 30\% cut in annual tile NO$_3$-N load would require treating tile flow 83\% of the time if targeting low flows but only 3\% if targeting high flows---a framework for sizing conservation practices to short windows of opportunity \parencite{Williams2018}. Managed-drainage monitoring reinforced this concentration: over a study period 287~mm of tile discharge (40.3\% of precipitation) occurred, 94\% of it when the outlet valve was open, and event discharge correlated strongly with the fraction of the event the valve was open (Pearson $r=0.86$) \parencite{Miller2022}.

\subsection{Travel times, water age, and groundwater interactions}

Modeling and tracer studies quantify how drainage shortens groundwater travel times and reorganizes the subsurface. In Iowa, groundwater travel-time distributions remained exponential across simulated drainage densities and incision depths; holding incision at 1.2~m, raising density from 0 to 0.0038~m$^{-1}$ shortened mean travel time and raised the tile contribution to baseflow to an upper bound of 37\% ($R^2=0.98$), while holding density constant and increasing incision from 0.3 to 2.7~m lowered mean travel time near-linearly (slope 1.45~yr~m$^{-1}$) and raised tile contribution from 30\% to 54\% \parencite{Schilling2012}. In Bear Creek, mean groundwater travel time fell from 5.57~yr at low drainage intensity to 1.07~yr at high intensity, far below the presettlement value of 181.6~yr, and intensive drainage short-circuited flow to riparian buffers---reducing land draining to buffers and cutting mean buffer residence time to 36.7~d or 15.5~d depending on the drainage-initiation scenario---thereby lowering buffer effectiveness for nitrate reduction \parencite{Schilling2015}. Analytical solutions confirm the mechanism: under tile drainage, groundwater levels are lower and decline faster, total groundwater discharge to the river increases, and the tile contribution to total discharge grows with tile density, though direct groundwater discharge to the river during recession is reduced because flow is diverted to the tiles \parencite{Liang2015}. Recession behavior of heavily tiled watersheds is correspondingly more linear than that of less-tiled watersheds \parencite{Schilling2008}, and differences in master recession curves were proposed as a diagnostic of tile contribution \parencite{Schilling2008}. Classical analytical treatments of transient recharge (Boussinesq solutions) provide the theoretical basis for such water-table and drain-flow dynamics \parencite{SINGH1980,Bouarfa2000}. Coexisting hydrograph-separation methods (SC-EMMA versus BFI-type baseflow programs) can nonetheless diverge substantially---bi-annual slowflow and baseflow indices differed by 5\% to 55\% across streams---underscoring method sensitivity when quantifying tile-influenced baseflow \parencite{Kronholm2015}.

\subsection{Controlled and managed drainage}

Controlled and managed drainage strongly modulate outflow volume and timing. In Sweden, controlled drainage reduced total drain outflow by 73--95\% relative to conventional drainage (e.g., 36, 14, and 120~mm versus 172, 240, and 450~mm over three years), concentrated more than 60\% of daily outflow in the 0--0.5~mm class, and produced lower retention constants and near-zero final storage; reduced storage was linked to higher peak flow, shorter lag, and shorter recession, effects that grew with groundwater level \parencite{Wesström2003}. Raising the drainage outlet significantly slowed water-table recession---by 29\% and 62\% in two site pairs---and lengthened the time for the water table to recede to 30 and 60~cm depths, whereas free-draining plots showed no such difference, confirming the effect was not seasonal \parencite{Saadat2017}. Agricultural drainage water management in Ohio was most effective at retaining water and reducing tile discharge during low-intensity events and dry antecedent conditions; the outlet valve was open on average 18.6\% of each three-day event yet passed 96\% of tile discharge, though average groundwater, soil moisture, and surface-flooding indicators did not differ significantly from free drainage \parencite{Miller2022}. In the Netherlands, subsurface water-infiltration systems significantly influenced water-table depth, reducing fluctuations at four of five sites and raising or deepening average water tables depending on applied water level and drainage base, with quadratic and delayed responses at several sites \parencite{vandeCraats2025}. Analyses of controlled-drainage and unreplicated watershed experiments highlight that method choice matters: paired-watershed, BACIP, beta/gamma-distribution, and time-series approaches can yield differing significance conclusions, with a beta distribution deemed most appropriate in one Minnesota dataset \parencite{Pagliari2021}. Cross-drain spacing further tunes recession, being most effective at 20~m spacing (reducing water-table recession by 13--32\%) and negligible by 40~m \parencite{Vaiciukynas2015}. Comparisons of drainage installation methods under Nordic conditions likewise revealed detectable performance differences: trenchless (T0) sections held groundwater levels 0.14--0.25~m higher than trench-installed (T1) sections in 60--90\% of observations, with the clearest differences when water levels rose above the 1.0~m drain depth and during autumn and winter \parencite{Salo2019}.

\subsection{Drain hydraulics, pressurization, and system design}

At the pipe scale, network pressurization introduces storage and delayed secondary peaks. In France, pressurization of the collector network during intense rainfall limited and sometimes reversed drainage discharge, temporarily storing infiltrated water as the water table rose and generating a four-stage sequence (first peak, backflow, second peak, recession) that increased the lag to the catchment outlet to 10--25~h \parencite{Henine2010}. A coupled groundwater--pipe model reproduced this behavior, with secondary peaks about 11--30~h after rainfall ended and estimated storage of 11--51~m$^3$ (38--46\% of stored recharge), the downstream network being most affected \parencite{Henine2014}. Drain flow rates were shown analytically to be governed by the system parameter $\sigma = K/(f^2 L^2)$, with maximum drain flow reaching 0.56~mm~h$^{-1}$ for $\sigma=10$~m$^{-1}$~h$^{-1}$ versus 0.025~mm~h$^{-1}$ for $\sigma=0.01$~m$^{-1}$~h$^{-1}$, and the constant-water-table-shape approximation improving as $\sigma$ increased \parencite{Bouarfa2000}. Combined tile/mole systems modeled with SEEP/W reproduced observed discharge well (e.g., $R^2=0.811$, index of agreement 0.94 in validation); a tile-only system was least effective, reducing drain discharge by 63\% and mean water-table depth by 72\% relative to the best combined systems, whose performance reflected the hydraulic conductivity above the mole channel \parencite{Tuohy2018}. Spatial patterns of drain flow were driven by topography and geology, with higher drain fractions in local depressions and amplified flow where hydraulic conductivity was higher \parencite{Mahmood2023}.

Component- and material-scale studies address inlets, media, and specialized drainage hardware relevant to flow and sediment. Surface-water-management structures such as tile risers and blind inlets altered flow duration and cumulative flow without necessarily changing discharge frequency, with direction of effect depending on the field, and blind-inlet infiltration declined from about 12--13 to 4--5~cm~h$^{-1}$ over six years, implying an 8--10-year service life under repeated tillage \parencite{Williams2020,Williams2019}. Replacing tile risers or open inlets with blind or gravel inlets substantially reduced sediment and phosphorus losses (e.g., TSS of 14.4 versus 40.4~kg~ha$^{-1}$~event$^{-1}$) \parencite{Feyereisen2015b}, and among alternative inlet designs a clear trade-off emerged between high flow capacity and sediment control \parencite{Li2017b}. In recycled concrete aggregate drainage media, permissible flow was strongly reduced by fines content (up to $\sim$73\% flow decline at 4\% fines), and filter-fabric opening size governed clogging potential \parencite{McCulloch2017}. Sheet-pipe subsurface drainage matched conventional pipe capacity and drained roughly 45\% more water with 55\% lower percolation, lowering water tables and improving aeration and rice yield \parencite{Tamura2021,Arif2021}. Numerical experiments on distribution pipes in saturated buffers showed two pipes could raise 10-day cumulative infiltration by up to $\sim$133\% relative to a single field-edge pipe, though at the cost of shorter subsurface travel distances \parencite{Jaynes2019b}. Deep-percolation and recharge estimation beneath drained lands, while uncertain, has long been recognized as part of the drainage water balance \parencite{BOOKMAN1959}.

\subsection{Coupled sediment, solute, and contaminant dynamics}

Because streamflow pathways carry material, hydrologic partitioning shapes sediment and solute export. In dual-drainage French fields, mean suspended-sediment yield was 0.79~t~ha$^{-1}$~yr$^{-1}$ with strong inter-annual variability (0.49 to 1.08~t~ha$^{-1}$~yr$^{-1}$); 89\% originated from surface drainage, and tile drainage contributed 46\% and 21\% of runoff but only 9\% and 11\% of sediment yield in the two years, with most export during the intense drainage period \parencite{Gaillot2023b,Grangeon2021}. A spatially distributed sediment-connectivity index captured surface source, pathway, and sink areas well but did not adequately represent connectivity where subsurface tile drainage controlled vertical flow \parencite{Calsamiglia2018}. Solute behavior in tile systems tended toward chemostatic transport, with antecedent conditions and storm size---not inlet density---dominating the response \parencite{Stops2022}. Tile drainage also delivered elevated concentrations of antibiotic-resistance genes relative to surface water, with concentrations correlated to daily flow rate \parencite{Rieke2018}. Converting surface runoff to subsurface drainage did not necessarily improve water quality: an infiltration-type raingarden retained 61\% of flow and 76--81\% of pollutants, whereas a drainage-type raingarden retained only 36\% of flow and could export more phosphorus and nitrate than it received, indicating that short-circuiting runoff through subsurface pipes can lower pollutant removal \parencite{Xu2022}.

\subsection{Specialized applications: salinity, bank stability, and valley-bottom drainage}

Beyond humid croplands, subsurface drainage governs streamflow-relevant processes in salinized, erosion-prone, and valley-bottom settings. In saline Chinese soils, subsurface drains lowered surface electrical conductivity over time and suppressed resalination better than undrained controls---daily resalination rates of 0.19--1.55\% under drainage versus 1.82\% for the control at one day after irrigation---while remaining economically favorable \parencite{Yan2013,Chang2012}. Passive drains lowering the water table stabilized gully banks in Ethiopia, cutting cumulative bank retreat by 52--84\% and soil loss by about 70\% at low cost, with a drained-bank groundwater table about 20\% lower than controls \parencite{Zegeye2021}. In a Burkina Faso valley bottom, under-phreatic permeability increased downstream from about $0.10\pm0.10$ to more than $20\pm10$~cm~h$^{-1}$, so no single infiltration process described the whole toposequence---an important consideration for efficient drainage design \parencite{Keita2014}.

\subsection{Modeling tile-drainage streamflow dynamics}

The synthesis rests on a spectrum of models whose consolidated performance is itself an established result. Equivalent-medium representations reproduced explicit tile simulations within 2\% of drained-volume targets and provided calibrated equivalent properties across tile densities for watershed-scale application \parencite{Thomas2016}. Physically based cold-region modeling captured tile timing, peak magnitude, and the seasonal absence of summer flow, improving when hourly steps and capillary rise were represented, and showed the capillary fringe strongly controlled drainable water and tile flow \parencite{Kompanizare2024}. A modified ODE hillslope formulation improved watershed-scale recession and flood-threshold timing specifically at artificially drained gauges, providing generalized seepage--storage relationships with physically interpretable thresholds \parencite{Fonley2021}. SWAT was rigorously tested against daily tile flow (NSE~$=0.5$, PBIAS~$=-1.4$\%) and combined with response analysis to assess basin-scale drainage impacts \parencite{Rahman2014}. A calibration-free time-expansion (travel-time) model matched calibrated SWAT skill in an intensely drained Indiana watershed using a single exponential distribution \parencite{Basu2010}, and multiple travel-time modeling approaches of differing sophistication converged closely in Bear Creek \parencite{Schilling2015}. Coupled SC-EMMA and hydrograph-recession frameworks using inexpensive sensors improved prediction of dissolved reactive phosphorus relative to any single method \parencite{Nazari2021}, and coupled groundwater--pipe models successfully represented pressurization, temporary retention, and flooded-outlet recovery beyond design capacity \parencite{Henine2014}. Collectively, these tools demonstrate that the streamflow signature of tile drainage---elevated low flows and baseflow, damped intermediate flows, modest large-flood effects, shortened and homogenized travel times, and scale-dependent flashiness and peak response---is both mechanistically explicable and reproducible across observational and modeling lines of evidence.

\printbibliography
\end{document}
